\documentclass[12pt]{article}
\usepackage[round]{natbib}
\usepackage[margin=2.5cm]{geometry}
\usepackage{setspace}
\onehalfspacing
\usepackage{graphicx} % Required for inserting images
\usepackage{amsthm}
\usepackage{amsmath,amssymb}
\let\dddot\relax     % relax AFTER amsmath is loaded
\let\ddddot\relax
\usepackage{accents}
\usepackage{amsfonts}
\usepackage{xfp}
\usepackage{comment}
\newcommand{\R}{\mathbb{R}}

\theoremstyle{definition}
\newtheorem{theorem}{Theorem}[section]
\newtheorem{proposition}{Proposition}[section]
\newtheorem{lemma}{Lemma}[section]
\newtheorem{definition}{Definition}[section]
\newtheorem{assumption}{Assumption}[section]

\newcommand{\ubar}[1]{\underaccent{\bar}{#1}}
% Auto-generated by pytask. Do not edit by hand.

\newcommand{\Pibaselineoneperiod}{0.225}
\newcommand{\Ubaselineoneperiod}{-2.100}
\newcommand{\tOnebaselineoneperiod}{0.275}
\newcommand{\tTwoAbaselineoneperiod}{4.505}
\newcommand{\tTwoBbaselineoneperiod}{-3.558}
\newcommand{\tTwoCbaselineoneperiod}{4.486}
\newcommand{\CAbaselineoneperiod}{-0.000}
\newcommand{\CBbaselineoneperiod}{-0.000}
\newcommand{\CCbaselineoneperiod}{-0.000}

\newcommand{\Pioneperiodlowoutsideoption}{0.225}
\newcommand{\Uoneperiodlowoutsideoption}{-2.100}
\newcommand{\tOneoneperiodlowoutsideoption}{0.275}
\newcommand{\tTwoAoneperiodlowoutsideoption}{4.505}
\newcommand{\tTwoBoneperiodlowoutsideoption}{-3.558}
\newcommand{\tTwoConeperiodlowoutsideoption}{4.486}
\newcommand{\CAoneperiodlowoutsideoption}{-0.000}
\newcommand{\CBoneperiodlowoutsideoption}{-0.000}
\newcommand{\CConeperiodlowoutsideoption}{-0.000}

\newcommand{\Pibaselinetwoperiod}{0.616}
\newcommand{\Ubaselinetwoperiod}{-4.241}
\newcommand{\tOnebaselinetwoperiod}{0.597}
\newcommand{\tTwoAbaselinetwoperiod}{1.025}
\newcommand{\tTwoBbaselinetwoperiod}{0.597}
\newcommand{\tTwoCbaselinetwoperiod}{-0.129}
\newcommand{\CAbaselinetwoperiod}{-0.300}
\newcommand{\CBbaselinetwoperiod}{0.168}
\newcommand{\CCbaselinetwoperiod}{0.144}

\newcommand{\Pitwoperiodlowoutsideoption}{0.621}
\newcommand{\Utwoperiodlowoutsideoption}{-4.372}
\newcommand{\tOnetwoperiodlowoutsideoption}{0.671}
\newcommand{\tTwoAtwoperiodlowoutsideoption}{1.473}
\newcommand{\tTwoBtwoperiodlowoutsideoption}{0.671}
\newcommand{\tTwoCtwoperiodlowoutsideoption}{-0.082}
\newcommand{\CAtwoperiodlowoutsideoption}{-0.300}
\newcommand{\CBtwoperiodlowoutsideoption}{0.131}
\newcommand{\CCtwoperiodlowoutsideoption}{0.165}

\newcommand{\Pitwoperiodimpatientprincipal}{0.660}
\newcommand{\Utwoperiodimpatientprincipal}{-4.032}
\newcommand{\tOnetwoperiodimpatientprincipal}{3.874}
\newcommand{\tTwoAtwoperiodimpatientprincipal}{2.986}
\newcommand{\tTwoBtwoperiodimpatientprincipal}{1.598}
\newcommand{\tTwoCtwoperiodimpatientprincipal}{-3.348}
\newcommand{\CAtwoperiodimpatientprincipal}{-0.030}
\newcommand{\CBtwoperiodimpatientprincipal}{-0.030}
\newcommand{\CCtwoperiodimpatientprincipal}{-0.173}


\title{A Two-Period Extension of Mattsson and Weibull's Principal--Agent Model}
\author{Orcan Ok}
\date{12 March 2026}

\begin{document}

\maketitle

\section{Introduction}

\cite{mattsson_weibull_2023} studies a single-period,
discrete-time principal–agent model with moral hazard and a finite set of outcomes. The agent can induce any probability distribution over outcomes, but doing so entails an effort cost that varies across distributions.

They focus on a class of effort-cost functions—\emph{Legendre-type} costs—that are invariant to merging equally paid outcomes. In particular,
if two outcomes yield the same wage to the agent, collapsing them into a single outcome leaves the agent’s optimal behavior unchanged.
For tractability, the analysis then specializes to the Kullback–Leibler (KL) divergence, a specific $f$-divergence within the broader class of Legendre-type costs.

These assumptions make the model analytically tractable and allow the authors to solve both the agent’s and the principal’s problems in closed
form. They then specialize to risk-neutral and constant-relative-risk-averse (CRRA) preferences to obtain interesting results, including the emergence of simple debt-like (threshold) optimal contracts.

In this write-up, I extend the model to a simple dynamic setting. Specifically, I consider two periods and allow the parties to sign a
long-term contract in which the second-period wage can be conditioned on the first-period outcome. The main question is whether the
debt-like (threshold) structure of optimal contracts survives in this environment. Intuitively, the prospect of future compensation
provides the principal with an additional incentive instrument to influence first-period effort. While this intertemporal linkage
could in principle complicate the contract, the extension remains tractable and yields a characterization that is closely related to the single-period case.

\section{The Original Model and Main Results in the Paper}
In the single-period baseline model, the principal hires an agent to perform a task. The agent’s effort is not verifiable and therefore not contractible; consequently, the contract can depend only on the realized output.

There is a finite set of outcomes, indexed by \(K:=\{1,\ldots,n\}\), and the set of probability distributions over \(K\) is $\Delta(K):=\Bigl\{p\in\mathbb{R}^{n}:\ \sum_{k=1}^{n}p_k=1\Bigr\}$.

The agent can induce any distribution \(p\in\Delta(K)\) at an effort cost \(\psi:\Delta(K)\to\mathbb{R}\), which is assumed to be continuous.

Each outcome \(k\in K\) is associated with a gross profit level \(\pi_k\in Q:=\{\pi_1,\ldots,\pi_n\}\subset\mathbb{R}\). 
Since the principal can verify only the realized output, she offers an output-contingent contract \(w:Q\to\mathbb{R}_{+}\), 
which specifies a nonnegative payment \(w_k:=w(\pi_k)\) for each \(\pi_k\in Q\). The restriction \(w_k\ge 0\) for all \(k\) captures limited liability of the agent.

The agent’s subutility from wage income is given by a function \(u:\mathbb{R}_{+}\to\mathbb{R}\), which is assumed to be continuous, strictly increasing, and unbounded above.

The agent’s effort cost is assumed to take the form of a Kullback--Leibler (KL) divergence:

\[
\psi(p)=\sum_{k\in K} p_k \ln\!\left(\frac{p_k}{p_k^{0}}\right),
\]
which is the maintained assumption in what follows.
The probabilities \(p^{0}=(p_k^{0})_{k=1}^{n}\in\Delta(K)\) can be interpreted as default probabilities, i.e.\ the outcome distribution that would obtain absent any effort.

In the KL case, the agent’s optimal choice \(p^{*}(w)\) takes the closed-form expression stated in Proposition~4 of the paper, where the agent's  optimal choice is
\[
p_k^{*}(w)=\frac{p_k^{0}e^{u(w_k)}}{\sum_{j\in K} p_j^{0}e^{u(w_j)}} \qquad \forall\,k\in K,
\]
and the associated achieved utility is
\[
U^{*}(w)=\ln\!\Bigl(\sum_{k\in K} p_k^{0}e^{u(w_k)}\Bigr).
\]

Solving the principal’s problem yields a canonical characterization of optimal contracts (Theorem~1 in Mattsson and Weibull). In particular, when \(\psi\) is of the KL form and \(u\) is
increasing and concave, an optimal contract can be written as a threshold--sharing rule:
\[
w^{*}(\pi_k)=s\!\left(\max\{0,\pi_k-t\}\right)\qquad \forall\,\pi_k\in Q,
\]
for some threshold \(t\in\mathbb{R}\) and some increasing function \(s:\mathbb{R}_{+}\to\mathbb{R}_{+}\) (the ``sharing rule'').

In particular, \(s\) satisfies the implicit equation
\begin{equation}\label{eq:s(x)}
s(x)=x-\frac{1}{u'(s(x))}+\frac{1}{u'(0)} \qquad \forall\,x\ge 0.
\end{equation}


\section{Extension: A Two-Period Contract}
\subsection{The Setting}

The environment is the same as in the baseline model, except that there are now two periods. At \(t=0\), the principal offers a long-term contract and can condition second-period wages on the first-period outcome.
I assume full commitment to the wage schedule: once the contract is signed, neither party can renegotiate after \(\pi_1\) is realized.\footnote{If the parties can renegotiate after \(\pi_1\) is realized, history-dependent
continuation terms are generally not credible: in period 2 there will typically exist a Pareto-improving renegotiation to the one-period optimal
contract for that period (given the relevant participation constraint). Anticipating this, the principal cannot credibly commit to future rewards or punishments, and the dynamic incentive effect disappears.
An alternative interpretation is that the principal has a commitment device that makes renegotiation infeasible—for example, wages may be pre-programmed and paid automatically in an irreversible way.}
A two-period contract is a pair \((w_1,w_2)\) where
\[
w_1:Q\to\mathbb{R}_+,\qquad w_2:Q\times Q\to\mathbb{R}_+,
\]
with the interpretation
\[
w_1(\pi_i)=:w_{1i},\qquad w_2(\pi_j\mid \pi_i)=:w_{2j\mid i}.
\]
Limited liability requires \(w_{1i}\ge 0\) and \(w_{2j\mid i}\ge 0\) for all \(i,j\in K\). The agent discounts future utility by \(\delta\in(0,1]\), while the principal discounts future profits by \(\beta\in(0,1]\).

The timing is as follows:
\begin{enumerate}
\item[\(t=0\):] The principal commits to a two-period contract \((w_1,w_2)\).
\item[\textbf{Period 1:}] The agent chooses \(p_1\in\Delta(K)\); an outcome \(\pi_1=\pi_i\) is realized; the wage \(w_{1i}\) is paid.
\item[\textbf{Period 2:}] After observing \(\pi_1=\pi_i\), the agent chooses \(p_2(\pi_i)\in\Delta(K)\); an outcome \(\pi_2=\pi_j\) is realized; the wage \(w_{2j\mid i}\) is paid.
\end{enumerate}

\subsection{Agent's Choice in Period 2}

Agent's problem in Period 2 is the same problem as in the one-period model, conditional on the history \(i\). Hence the optimal choice again takes the form:
\begin{equation}\label{eq:p2j}
p_{2j}^*(i)
=
\frac{p_j^0 e^{u(w_{2j\mid i})}}{\sum_{\ell\in K}p_\ell^0 e^{u(w_{2\ell\mid i})}},
\qquad \forall\, j\in K.
\end{equation}

The agent’s continuation utility at the beginning of period~2 is therefore
\begin{equation}\label{eq:u2_agent}
U_2^*(i)
=
\max_{p_2\in\Delta(K)}\left\{\sum_{j\in K} p_{2j}\,u(w_{2j\mid i})-\psi(p_2)\right\}
=
\ln\!\left(\sum_{\ell\in K}p_\ell^0 e^{u(w_{2\ell\mid i})}\right).
\end{equation}

\subsection{Agent's Choice in Period 1}

Solving the agent’s problem in Period 1 by standard Lagrangian methods yields
\begin{equation}\label{eq:p2i}
p_{1i}^*
=
\frac{p_i^0\,e^{u(w_{1i})}\,Z_i^{\delta}}
{\sum_{m\in K}p_m^0\,e^{u(w_{1m})}\,Z_m^{\delta}}
=
\frac{p_i^0\,e^{u(w_{1i})+\delta U_2^*(i)}}
{\sum_{m\in K}p_m^0\,e^{u(w_{1m})+\delta U_2^*(m)}},
\qquad \forall\, i\in K.
\end{equation}
The associated ex-ante utility at the beginning of period~1 is
\begin{equation}\label{eq:u1_agent}
U^*(w_1,w_2)
=
\ln\!\left(\sum_{i\in K}p_i^0\,e^{u(w_{1i})}Z_i^{\delta}\right)
=
\ln\!\left(\sum_{i\in K}p_i^0\,e^{u(w_{1i})+\delta U_2^*(i)}\right).
\end{equation}

Relative to the one-period model (Proposition~4), the only difference is that the agent’s effective subutility for outcome \(i\) is \(u(w_{1i})+\delta U_2^*(i)\), i.e.\ current 
utility plus the discounted continuation value induced by the history-dependent second-period contract.

\subsection{Optimal Contract Chosen by the Principal}

Solving the principal's problem using Lagrangian method gives the following theorem:

\begin{theorem}\label{thm:two_period_structure}
Assume that the agent’s wage subutility \(u:\mathbb{R}_+\to\mathbb{R}\) is twice differentiable with \(u'>0\) and \(u''\le 0\), and that the effort cost is of the KL form. 
Then the optimal two-period contract admits a threshold--sharing representation. In particular, there exist thresholds \(t_1\in\mathbb{R}\) and \(t_2:K\to\mathbb{R}\) such that
\begin{equation}\label{eq:w1_general}
w_1^*(\pi_i)=s\!\left(\max\{0,\pi_i+C_i-t_1\}\right),
\qquad \forall\, i\in K,
\end{equation}
and
\begin{equation}\label{eq:w2_general}
w_2^*(\pi_j\mid \pi_i)=s\!\left(\max\{0,\pi_j-t_2(i)\}\right),
\qquad \forall\, i,j\in K.
\end{equation}
Here \(s:\mathbb{R}_+\to\mathbb{R}_+\) is the sharing rule characterized by \eqref{eq:s(x)}, and
\[
C_i=\beta\sum_{j\in K}p_{2j}^*(i)\,(\pi_j-w_{2j\mid i})
\]
denotes the principal’s discounted continuation profit following history \(i\).
\end{theorem}
\begin{proof}
See Appendix.
\end{proof}
Thus, in period~1, the principal behaves as if there were a one-period contract and the relevant ``effective output'' were \(\pi_i+C_i\), i.e.\ current output plus discounted expected continuation profit.
Not surprisingly, in period~2, the principal faces a one-period problem given the first-period outcome \(i\), hence the second-period contract is the same as it would be in a standalone one-period setting.

\subsubsection{Numerics and Discussion}
It is worth emphasizing that, although the two-period problem---especially the principal’s problem (see Appendix)---appears substantially more complicated than its 
one-period counterpart, the resulting optimal behavior remains closely linked to the single-period solution. In period~1, the agent behaves as if facing a one-period 
problem with an effective utility equal to current utility plus discounted expected utility tomorrow. Likewise, the principal behaves as if facing a one-period 
problem in which the effective output in history \(i\) is \(\pi_i\) plus the discounted expected profit induced by the period-2 contract.

It is natural to complement the analytical characterization with a numerical comparison between (i) an environment in which the parties can sign a joint two-period 
contract and (ii) a benchmark in which contracting is restricted to one-period agreements in each period.

Since the principal can always replicate the one-period benchmark within the two-period contract space—by choosing continuation terms that do not depend on the 
first-period outcome (i.e.\ setting \(t_2(i)\) constant across \(i\))—the option to write a joint two-period contract can only weakly increase the principal’s expected profit.

The welfare implications for the agent are less clear. In particular, the principal may use history dependence to create ``career concerns'': poor 
first-period performance can be followed by an unfavorable continuation contract, whereas good performance can be rewarded with more favorable terms. 
Such intertemporal incentives may induce the agent to exert more effort in period~1 than under a stand-alone one-period contract. As a result, 
the agent may be worse off relative to the benchmark of two separate one-period contracts, even if participation holds ex ante.

At the same time, the opposite pattern is also possible. If the principal is relatively impatient (i.e.\ \(\beta\) is small), she may prefer to 
backload compensation and use generous promised continuation terms to strengthen first-period incentives at relatively low present-value cost. 
In that case, a joint two-period contract can increase both expected profit and expected wage payments, and the agent may end up strictly better off than under repeated one-period contracting.

I begin with the numerical example in the appendix of Mattsson and Weibull.\footnote{For the code implementing the numerics from here on, see https://github.com/orcanok/mattsson-weibull-extension}
There are three outcomes \(K=\{1,2,3\}\) with gross profits \(\pi\in\{-1,0,1\}\) 
and baseline probabilities \(p^0=(0.5,0.3,0.2)\). The agent has unit relative risk aversion (\(\rho=1\)), external income \(\varepsilon=0.05\), 
and a one-period outside option \(\bar u^{(1)}=\ln(0.12)\approx -2.12\).\footnote{Since utility is measured on an arbitrary cardinal scale 
(and, for \(\rho=1\), \(u(w)=\theta\ln(w+\varepsilon)\)), negative reservation utility levels are not problematic.} In this case, the one-period 
optimal contract is a threshold contract with \(t\approx \tOnebaselineoneperiod\), yielding an expected profit of approximately \(\Pibaselineoneperiod\) 
for the principal while giving the agent expected utility \(\fpeval{round(\Ubaselineoneperiod, 5)}\), i.e. the participation constraint almost binds.

Assume both parties are fully patient, i.e.\ \(\beta=\delta=1\). The agent’s outside option for the \emph{two-period} relationship is
\[
\bar u \;=\; 2\ln(0.12)\;\approx\;-4.241 .
\]
If the parties contract separately in each period using the one-period optimal contract (with the same per-period outside option), 
the principal’s expected total profit is approximately \(\fpeval{round(2*\Pibaselineoneperiod, 5)}\), while the agent’s expected 
total utility is approximately \(\fpeval{round(2*\Ubaselineoneperiod, 5)}\) (i.e.\ almost her outside option). If instead the parties can sign a 
long-term two-period contract, the principal optimally chooses a history-dependent wage schedule of the form
\[
w_1^*(\pi_i)=s\!\left(\max\{0,\pi_i+C_i-t_1\}\right),\qquad \forall\, i\in K,
\]
with \(t_1\approx \tOnebaselinetwoperiod\) and continuation terms
\[
C_1 \approx \CAbaselinetwoperiod,\qquad C_2 \approx \CBbaselinetwoperiod,\qquad C_3 \approx \CCbaselinetwoperiod.
\]
Continuation wages take the analogous form
\[
w_2^*(\pi_j\mid \pi_i)=s\!\left(\max\{0,\pi_j-t_2(i)\}\right),\qquad \forall\, i,j\in K,
\]
where
\[
t_2(1)\approx \tTwoAbaselinetwoperiod,\quad t_2(2)\approx \tTwoBbaselinetwoperiod,\quad t_2(3)\approx \tTwoCbaselinetwoperiod.
\]
Under this joint contract, the principal’s expected profit rises to about \(\Pibaselinetwoperiod\), while the agent’s expected utility remains at \(\Ubaselinetwoperiod\); 
thus the participation constraint binds and the agent is held at her outside option. As expected, allowing a joint two-period contract makes the principal strictly better off in this example.

Next, lower the agent’s \emph{one-period} outside option to \(\ln(0.05)\approx -2.996\). As in the original paper, the one-period optimal 
contract remains essentially unchanged (with \(t\approx \tOneoneperiodlowoutsideoption\)), so the induced allocation and the principal’s profit are unchanged; 
however, participation is now slack and the agent receives strictly more than her reservation utility (e.g.\ \(\Uoneperiodlowoutsideoption > -2.996\) 
in the one-period benchmark). If, on the other hand, the parties can sign a two-period contract, the principal again chooses a contract of the form before. The optimal parameters are approximately
\[
t_1 \approx \tOnetwoperiodlowoutsideoption,\qquad
t_2(1)\approx \tTwoAtwoperiodlowoutsideoption,\quad
t_2(2)\approx \tTwoBtwoperiodlowoutsideoption,\quad
t_2(3)\approx -0.08,
\]
\[
C_1 \approx \CAtwoperiodlowoutsideoption,\qquad
C_2 \approx \CBtwoperiodlowoutsideoption,\qquad
C_3 \approx \CCtwoperiodlowoutsideoption.
\]
Under this joint contract, the principal’s expected profit is about \(\Pitwoperiodlowoutsideoption\), while the agent’s expected utility is \(\Utwoperiodlowoutsideoption\), 
which is lower than under repeated one-period contracting \(\fpeval{round(2*\Uoneperiodlowoutsideoption, 3)}\). Thus, allowing a joint two-period contract can increase the principal’s profit at the expense of the agent.

To illustrate that joint contracting can also benefit the agent, revert to the baseline one-period outside option \(\bar u^{(1)}\approx -2.12\) and set \(\delta=1\) 
so that \(\bar u^{(2)}\approx -4.241\).\footnote{With \(\delta=1\), the two-period outside option equals \(\bar u^{(2)}=(1+\delta)\bar u^{(1)}\approx -4.241\).} 
Suppose further that the principal is very impatient, with \(\beta=0.1\), while all other parameters remain unchanged. The optimal two-period contract then has approximately
\[
t_1 \approx \tOnetwoperiodimpatientprincipal,\qquad
t_2(1)\approx \tTwoAtwoperiodimpatientprincipal,\quad
t_2(2)\approx \tTwoBtwoperiodimpatientprincipal,\quad
t_2(3)\approx \tTwoCtwoperiodimpatientprincipal,
\]
\[
C_1 \approx \CAtwoperiodimpatientprincipal,\qquad
C_2 \approx \CBtwoperiodimpatientprincipal,\qquad
C_3 \approx \CCtwoperiodimpatientprincipal.
\]
It is interesting to note that $t_1$ is so high that principal gets $0$ wage in period-1 regardless of the output. In this case, the principal’s 
expected profit is about \(\Pitwoperiodimpatientprincipal\), and the agent’s expected utility is \(\Utwoperiodimpatientprincipal\), which is higher than under repeated one-period 
contracting \(\fpeval{round(2*\Uoneperiodlowoutsideoption, 3)}\). 
Hence, the possibility of a joint two-period contract can also generate a Pareto improvement.
\section{Conclusion}

This write-up extends Mattsson and Weibull’s analytically solvable principal--agent model to a two-period environment in which the principal can 
condition continuation wages on first-period performance. The extension preserves the structure of the agent’s induced distributions: period-2 
behavior is identical to the one-period problem conditional on the first-period history, and period-1 behavior depends on an effective utility 
that aggregates current utility and discounted continuation value.

Analytically, the optimal wage schedules retain a simple canonical form. In period~2, conditional on the realized first-period outcome, 
the principal faces a one-period problem, so the continuation wage schedule has the same threshold--sharing structure as in the baseline model. 
In period~1, the wage schedule again has a threshold--sharing form, but it is written on a continuation-adjusted performance measure: 
current output plus the principal’s discounted expected continuation profit \(C_i\). 

Numerically, allowing a joint two-period contract can strictly increase the principal’s expected profit relative to repeated one-period contracting. 
The welfare implications for the agent are ambiguous: history-dependent continuation terms may discipline first-period behavior and reduce the agent’s 
expected utility relative to the benchmark, but when the principal is sufficiently impatient, backloading incentives through the continuation 
contract can also raise the agent’s expected utility, generating a Pareto improvement. A natural direction for future work is to explore whether similar qualitative patterns persist in longer horizons.

\section{Appendix: Proof for Theorem~\ref{thm:two_period_structure}}
\label{app:general_u}
Let $D:=\sum_{i\in K}p_i^0\,e^{u(w_{1i})}\,Z_i^\delta$. The principal's expected profit is:

\begin{equation}\label{eq:principal_problem}
  \Pi^*(w_1,w_2)
=
\sum_{i\in K}
\frac{p_i^0\,e^{u(w_{1i})}\,Z_i^\delta}{D}
\Bigg[
(\pi_i-w_{1i})
+\beta\cdot \frac{1}{Z_i}\sum_{j\in K}p_j^0 e^{u(w_{2j\mid i})}\,(\pi_j-w_{2j\mid i})
\Bigg].
\end{equation}


The principal chooses $(w_1,w_2)$ to maximize $\Pi^*(w_1,w_2)$ subject to participation and limited liability. Assume that the agent's outside 
option is $\bar{u}$. The participation constraint of the agent becomes: $U^*(w_1,w_2)=\ln D\ge \bar u\quad\Longleftrightarrow\quad D\ge e^{\bar u}.$


Let \(\eta\ge 0\) be the multiplier on participation and let \(\mu_{1i}\ge 0\) and
\(\mu_{2j\mid i}\ge 0\) be the multipliers on limited liability. The Lagrangian is
\[
\mathcal{L}
=
\Pi^*(w_1,w_2)
+\eta\big(\ln D-\bar u\big)
+\sum_{i\in K}\mu_{1i}w_{1i}
+\sum_{i\in K}\sum_{j\in K}\mu_{2j\mid i}w_{2j\mid i}.
\]
Define the principal's continuation (net) profit given history \(i\) as
\[
\Pi_i
:=
(\pi_i-w_{1i})
+\beta\sum_{j\in K}p_{2j}^*(i)\,(\pi_j-w_{2j\mid i}),
\qquad
\Pi^*:=\sum_{i\in K}p_{1i}^*\Pi_i.
\]
Then, the Karush--Kuhn--Tucker (KKT) condition for \(w_{1i}\) is
\begin{equation}\label{eq:kkt_1_app}
0
=
\frac{\partial \mathcal{L}}{\partial w_{1i}}
=
-\,p_{1i}^*
+
p_{1i}^*u'(w_{1i})\big(\Pi_i-\Pi^*\big)
+\eta\,p_{1i}^*u'(w_{1i})
+\mu_{1i}.
\end{equation}
Define conditional period-2 net profit in state \(j\) after history \(i\) by
\[
\tilde\Pi_{ij}:=\pi_j-w_{2j\mid i},
\qquad
\tilde\Pi_i:=\sum_{\ell\in K}p_{2\ell}^*(i)\tilde\Pi_{i\ell}.
\]
Then, the KKT condition for \(w_{2j\mid i}\) is
\begin{equation}\label{eq:kkt_2_app}
-\beta
+
\beta\,u'(w_{2j\mid i})\big(\tilde\Pi_{ij}-\tilde\Pi_i\big)
+
\delta\,u'(w_{2j\mid i})\big(\Pi_i-\Pi^*+\eta\big)
+
\frac{\mu_{2j\mid i}}{p_{1i}^*p_{2j}^*(i)}
=0.
\end{equation}
Assume \(u:\mathbb{R}_+\to\mathbb{R}\) is twice differentiable with \(u'>0\) and \(u''\le 0\).
Define
\begin{equation}\label{eq:g_def_general}
g(w):=w+\frac{1}{u'(w)},\qquad w\ge 0.
\end{equation}
Since $g'(w) \ge 0$, the function \(g\) is strictly increasing on \(\mathbb{R}_+\) and hence invertible. Define the sharing rule \(s:\mathbb{R}_+\to\mathbb{R}_+\) by $s(x):=g^{-1}\!\big(g(0)+x\big)$ , $x\ge 0.$
Equivalently, \(s\) is characterized implicitly by
\begin{equation}\label{eq:s_general_def_implicit}
x=s(x)+\frac{1}{u'(s(x))}-\frac{1}{u'(0)},\qquad \forall x\ge 0,
\end{equation}
which is the same condition as \eqref{eq:s(x)} in Theorem~1.
Define
\[
A_i:=\delta(\Pi_i-\Pi^*+\eta),
\qquad
\Gamma_i:=\frac{A_i}{\beta}-\tilde\Pi_i.
\]
We get
\begin{equation}\label{eq:period2_g_identity_general}
g(w_{2j\mid i})
=w_{2j\mid i}+\frac{1}{u'(w_{2j\mid i})}
=\pi_j+\Gamma_i.
\end{equation}
Since \(g\) is strictly increasing, the interior solution is \(w_{2j\mid i}=g^{-1}(\pi_j+\Gamma_i)\).
Imposing limited liability \(w_{2j\mid i}\ge 0\), define the history-dependent threshold
\begin{equation}\label{eq:t2_def_general}
t_2(i):=g(0)-\Gamma_i.
\end{equation}
If \(\pi_j\le t_2(i)\), then \(\pi_j+\Gamma_i\le g(0)\) and the limited-liability corner implies \(w_{2j\mid i}=0\). If \(\pi_j>t_2(i)\), then
\(\pi_j+\Gamma_i=g(0)+(\pi_j-t_2(i))\) and thus
\[
w_{2j\mid i}=g^{-1}\!\big(g(0)+(\pi_j-t_2(i))\big)=s(\pi_j-t_2(i)).
\]
Combining the corner and interior regions yields
\begin{equation}\label{eq:w2_general_derived}
w_2^*(\pi_j\mid \pi_i)=s\!\left(\max\{0,\pi_j-t_2(i)\}\right),
\qquad \forall\, i,j\in K.
\end{equation}
Now consider \eqref{eq:kkt_1_app} in the interior region where \(w_{1i}>0\) so \(\mu_{1i}=0\). Some algebra gives
\[
\frac{1}{u'(w_{1i})}=\pi_i-w_{1i}+C_i-\Pi^*+\eta
\quad\Longleftrightarrow\quad
g(w_{1i})=\pi_i+C_i-\Pi^*+\eta.
\]
Thus the interior solution is \(w_{1i}=g^{-1}(\pi_i+C_i-\Pi^*+\eta)\). Imposing limited liability \(w_{1i}\ge 0\), define
\begin{equation}\label{eq:t1_def_general}
t_1:=g(0)+\Pi^*-\eta.
\end{equation}
If \(\pi_i+C_i\le t_1\), then \(w_{1i}=0\); otherwise,
\[
w_{1i}=g^{-1}\!\big(g(0)+(\pi_i+C_i-t_1)\big)=s(\pi_i+C_i-t_1).
\]
Hence
\begin{equation}\label{eq:w1_general_derived}
w_1^*(\pi_i)=s\!\left(\max\{0,\pi_i+C_i-t_1\}\right),
\qquad \forall\, i\in K.
\end{equation}


\bibliographystyle{apalike} 
\bibliography{references}

\end{document}
